% --- Codifica e Lingua ---
\usepackage[utf8]{inputenc}
\usepackage[T1]{fontenc}
\usepackage[italian]{babel}

% --- Pacchetti base e Layout ---
\usepackage{amsmath, amssymb}
\usepackage{geometry}
\geometry{a4paper, margin=1.75cm}
\usepackage{titlesec}
\usepackage[colorlinks=true, linkcolor=blue, urlcolor=cyan]{hyperref}

% --- Setup Grafica e TikZ ---
\usepackage{tikz}
\usetikzlibrary{positioning, arrows.meta, shapes.geometric, shapes.multipart}

\tikzset{
    entity/.style={rectangle, draw=black, thick, fill=blue!5, minimum width=2.5cm, minimum height=1cm, align=center},
    weak_entity/.style={rectangle, draw=black, thick, double, double distance=2pt, fill=blue!5, minimum width=2.5cm, minimum height=1cm, align=center},
    relationship/.style={diamond, draw=black, thick, fill=green!5, minimum width=2.2cm, minimum height=1cm, aspect=1.5, align=center},
    ident_relationship/.style={diamond, draw=black, thick, double, double distance=2pt, fill=green!5, minimum width=2.2cm, minimum height=1cm, aspect=1.5, align=center},
    attribute/.style={ellipse, draw=black, thin, fill=yellow!5, minimum width=1.8cm, minimum height=0.8cm, align=center, font=\small},
    id_attribute/.style={ellipse, draw=black, thick, fill=yellow!15, minimum width=1.8cm, minimum height=0.8cm, align=center, font=\small\bfseries},
    isa/.style={isosceles triangle, draw=black, thick, fill=gray!20, minimum width=1cm, minimum height=0.8cm, shape border rotate=90, align=center},
    link/.style={draw, black, thick},
    table_schema/.style={rectangle split, rectangle split parts=#1, rectangle split horizontal, draw=black, thick, fill=gray!10, minimum height=0.7cm, align=center, font=\small\ttfamily},
    fk_arrow/.style={-{Stealth[scale=1.2]}, thick, color=blue!70!black, rounded corners=4pt}
}

% --- Configurazione Testo ---
\setlength{\parindent}{0pt}
\setlength{\parskip}{0.3\baselineskip}
\titleformat{\chapter}[hang]{\normalfont\huge\bfseries}{\thechapter.}{20pt}{\Huge}
\titlespacing*{\chapter}{0pt}{-20pt}{30pt}

% --- Configurazione SQL (listings) ---
\usepackage{listings}
\usepackage{xcolor}
% Inserisci queste righe in param.tex, subito dopo \usepackage{xcolor}
\definecolor{headerblue}{HTML}{00CCFF}
\definecolor{rowblue}{HTML}{E0FFFF}
\definecolor{textdarkblue}{HTML}{000080}
\lstdefinelanguage{sql_generic}{
    morekeywords={SELECT, FROM, WHERE, GROUP, BY, HAVING, JOIN, LEFT, RIGHT, FULL, INNER, OUTER, UNION, INTERSECT, EXCEPT, CREATE, TABLE, ALTER, DROP, INSERT, UPDATE, DELETE, PRIMARY, FOREIGN, KEY, REFERENCES, CONSTRAINT, ON, AND, OR, NOT, IN, EXISTS, DISTINCT, AS, AVG, COUNT, MIN, MAX, SUM, DECIMAL, CHAR, VARCHAR, INT, DATE, DEFAULT, CHECK, CASCADE, RESTRICT, SCHEMA, VIEW},
    sensitive=false,
    morecomment=[l]{--},
    morestring=[b]',
}
\lstset{
    language=sql_generic,
    basicstyle=\ttfamily\small,
    keywordstyle=\color{blue}\bfseries,
    commentstyle=\color{gray},
    stringstyle=\color{red},
    frame=single,
    breaklines=true,
    captionpos=b
}

% --- Altri Pacchetti ---
\usepackage{colortbl}
\usepackage{array}
\newcolumntype{C}[1]{>{\centering\arraybackslash}p{#1}}

